% Part: normal-modal-logic
% Chapter: frame-correspondence
% Section: frames

\documentclass[../../../include/open-logic-section]{subfiles}

\begin{document}

\olfileid[hi]{nml}{frd}{fra}

\olsection{फ़्रेम}

\begin{defn}
  \emph{फ़्रेम} एक युग्म $\mModel{F} = \tuple{W,R}$ है, जहाँ $W$ जगतों का
  अरिक्त समुच्चय और $R$, $W$ पर द्विपदी संबंध है। कोई प्रतिरूप
  $\mModel{M}$ फ़्रेम $\mModel{F} = \tuple{W,R}$ पर \emph{आधारित} है, तभी
  जब किसी मूल्यांकन $V$ के लिए $\mModel{M} = \tuple{W, R, V}$।
\end{defn}

\begin{defn}
  यदि $\mModel{F}$ कोई फ़्रेम है, तो हम कहते हैं कि $!A$, $\mModel{F}$ में
  \emph{मान्य है}, $\mModel{F} \Entails !A$, यदि $\mModel{F}$ पर आधारित
  प्रत्येक प्रतिरूप~$\mModel{M}$ के लिए $\mSat{M}{!A}$।
  
  यदि $\mClass{F}$ फ़्रेमों का वर्ग है, तो हम कहते हैं कि $!A$,
  $\mClass{F}$ में \emph{मान्य है}, $\mClass{F} \Entails !A$, तभी जब
  प्रत्येक फ़्रेम~$\mModel{F} \in \mClass{F}$ के लिए
  $\mModel{F} \Entails !A$।
\end{defn}

फ़्रेम इसलिए रोचक हैं कि स्कीमाओं और अभिगम्यता संबंध~$R$ के गुणों के बीच
अनुरूपता फ़्रेमों के स्तर पर होती है, \emph{प्रतिरूपों के स्तर पर नहीं}।
उदाहरणार्थ, यद्यपि \Ax{T} सभी स्वपरावर्ती प्रतिरूपों में सत्य है, प्रत्येक
ऐसा प्रतिरूप जिसमें \Ax{T} सत्य है स्वपरावर्ती नहीं होता। किंतु यह
\emph{सत्य} है कि न केवल \Ax{T} सभी स्वपरावर्ती \emph{फ़्रेमों} पर
\emph{मान्य} है, बल्कि प्रत्येक ऐसा फ़्रेम जिसमें \Ax{T} मान्य है,
स्वपरावर्ती भी है।

\begin{rem}
फ़्रेमों के किसी वर्ग में मान्यता, प्रतिरूपों के किसी वर्ग में मान्यता की
धारणा का विशेष मामला है: $\mClass{F} \Entails !A$ तभी जब
$\mClass{C} \Entails !A$, जहाँ $\mClass{C}$, $\mClass{F}$ के किसी फ़्रेम
पर आधारित सभी प्रतिरूपों का वर्ग है।

स्पष्ट है कि यदि कोई !!{formula} या स्कीमा मान्य है, अर्थात् \emph{सभी}
प्रतिरूपों के वर्ग के सापेक्ष मान्य है, तो वह फ़्रेमों के किसी भी वर्ग
$\mClass{F}$ के सापेक्ष भी मान्य है।
\end{rem}

\end{document}
